\scanpage{147}
\mspage{63}

\[
\left\{\begin{aligned}
(f^{M}_{S})^{*}&: S^{\wedge}\longrightarrow M\\
(f^{N}_{T})^{\unsure{\prime|*}}&: T^{\wedge}\longrightarrow N,
\end{aligned}\right.
\tag*{(4.2.12)}
\]

à savoir les prolongements canoniques (\unsure{qui commutent|par commutation} aux petites $\unsure{\underrightarrow{\lim}|\underline{\lim}}$) des foncteurs d'inclusion $S\lhook\joinrel\longrightarrow M$, $T\lhook\joinrel\longrightarrow N$.

La question essentielle, à présent, c'est évidemment de construire un \deleted{\unsure{$f$}} adjoint à gauche pour le foncteur (4.2.9). Il revient au $=$, pour $(s,t)$ dans $S\times T$, de définir un objet

\[
s\boxtimes t\in\operatorname{Ob} M\boxtimes N
\tag*{(4.\deleted{\unsure{1}}2.13)}
\]

tel que l'on ait, pour $F$ variable dans $M\boxtimes N$, un isom. fonctoriel en $F$

\[
\underbrace{\unsure{\rho|\varrho}(F)(s,t)}_{\substack{\Vert\,\mathrm{def}\\ \deleted{\unsure{\mathrm{def}}}F(s,t)}}\simeq\operatorname{Hom}(s\boxtimes t,F):
\]

Donc il faut prouver ceci

% slanted
\margin{Pas prouvé, peut-être faux !}

\uline{Lemme} 4.2.3. Soient $M,N$ deux catégories pseudo-topos, et $\deleted{\unsure{x}}s\in\operatorname{Ob}M$, $\deleted{\unsure{y}}t\in\operatorname{Ob}N$, et considérons sur $M\boxtimes N\overset{\mathrm{def}}{=}\underline{\operatorname{Hom}}^{!!}(M^{\mathrm{o}}\times N^{\mathrm{o}},\mathrm{Ens})$ le foncteur $F\longmapsto F(\deleted{\unsure{x}}s,\deleted{\unsure{y}}t)$. Ce foncteur est \uline{représentable} (par un objet que je note $\deleted{\unsure{x}}s\boxtimes\deleted{\unsure{y}}t$).

L'idée naturelle serait de définir $\deleted{\unsure{x}}s\boxtimes\deleted{\unsure{y}}t$ par

\[
(\deleted{\unsure{x}}s\boxtimes\deleted{\unsure{y}}t)(\deleted{\unsure{s}}x,y)
=\deleted{\unsure{x}}s(\deleted{\unsure{s}}x)\times\deleted{\unsure{y}}t(y)
\unsure{\equiv|=}\operatorname{Hom}_{M}(x,s)\times\operatorname{Hom}_{N}(y,t)
\]
% archive sheet number: 146
